\documentclass{article}

\usepackage{amsmath,amssymb}
\usepackage{xurl}
\usepackage[hidelinks]{hyperref}
\setlength{\emergencystretch}{2em}

% Shared commands for notes in docs/paper-gaps/.

\DeclareUnicodeCharacter{2080}{\ifmmode{}_{0}\else\textsubscript{0}\fi}
\DeclareUnicodeCharacter{2081}{\ifmmode{}_{1}\else\textsubscript{1}\fi}
\DeclareUnicodeCharacter{2082}{\ifmmode{}_{2}\else\textsubscript{2}\fi}
\DeclareUnicodeCharacter{2099}{\ifmmode{}_{n}\else\textsubscript{n}\fi}

\newcommand{\Lean}{\textsc{Lean}}
\ExplSyntaxOn
\NewDocumentCommand{\leanid}{m}{
  \ifmmode
    \textsf{\detokenize{#1}}
  \else
    \group_begin:
    \urlstyle{sf}
    \tl_set:Nn \l_tmpa_tl {#1}
    \tl_replace_all:Nnn \l_tmpa_tl {\_} {_}
    \exp_args:NV \path \l_tmpa_tl
    \group_end:
  \fi
}
\ExplSyntaxOff
\providecommand{\tr}{\operatorname{tr}}
\newcommand{\tnleanIssueBase}{https://github.com/LionSR/TNLean/issues/}
\newcommand{\tnleanPrBase}{https://github.com/LionSR/TNLean/pull/}
\newcommand{\tnleanDocBase}{https://sirui-lu.com/TNLean/docs/}
\newcommand{\ghissue}[1]{\href{\tnleanIssueBase#1}{GitHub issue \##1}}
\newcommand{\ghpr}[1]{\href{\tnleanPrBase#1}{PR \##1}}
\newcommand{\doclink}[1]{\href{\tnleanDocBase#1.html}{\path{#1}}}

% Dirac notation and the identity operator, as in blueprint/src/macros/common.tex.
% \providecommand keeps note-local definitions authoritative.
\usepackage{dsfont}
\providecommand{\ket}[1]{|#1\rangle}
\providecommand{\bra}[1]{\langle #1|}
\providecommand{\braket}[2]{\langle #1 | #2 \rangle}
\providecommand{\ketbra}[2]{|#1\rangle\!\langle #2|}
\providecommand{\Id}{\mathds{1}}
\providecommand{\one}{\mathds{1}}
\providecommand{\C}{\mathbb{C}}
\providecommand{\MN}[1]{M_{#1}(\C)}
\providecommand{\MD}{M_D(\C)}

% Verdict marker: \gapnote{<kind>}{<status>}. Kind is the policy
% classification (clarification, local-correction, scope-restriction,
% unfaithful, false-source, open-gap); status is open, wip (elimination
% underway), resolved, or historical. Machine-read by the site index;
% prints a verdict line.
\newcommand{\gapnote}[2]{\par\noindent\textbf{Verdict:} #1 (#2).\par}


% Fallbacks keep this note compilable when a different command.tex is found first.
\providecommand{\leanid}[1]{\textsf{\detokenize{#1}}}
\providecommand{\gapnote}[2]{\par\noindent\textbf{Verdict:} #1 (#2).\par}

\title{The Constants in the Entropy Gain of Hastings' Area Law}
\author{}
\date{2026-10-22}

\begin{document}
\maketitle
\gapnote{clarification}{resolved}

\section{The source passage}

Hastings' proof of the one-dimensional area law~\cite{Hastings2007AreaLaw}
(arXiv:0705.2024v4) reduces the strengthened doubling inequality
(lines~201--204) to the bound
\begin{equation}\label{eq:relent}
  S(\rho_{j-l+1,j})+S(\rho_{j+1,j+l})-S(\rho_{j-l+1,j+l})
  \ge\bigl(1-2\epsilon\bigr)\frac{l}{\xi'}-\ln C_1-C_2,
  \qquad\epsilon=C_1e^{-l/\xi'},
\end{equation}
for adjacent intervals of $l$ sites (lines~238--242). Lines~413--503 prove it.
The operators $O_B,O_L,O_R$ of Lemma~1 satisfy
$\lVert O_BO_LO_R-P_0\rVert\le\epsilon$. Let $\rho$ be the ground-state
marginal on the $2l$ sites, $\sigma=\rho_{1,j}\otimes\rho_{j+1,N}$,
$x=\operatorname{tr}(O_B\sigma)$ and $y=\operatorname{tr}(O_LO_R\sigma)$.
The proof shows $\operatorname{tr}(O_B\rho)\ge1-2\epsilon$, bounds $x$
(eq.~\texttt{xbd}), and bounds the mutual information below by the binary
relative entropy
\begin{equation}\label{eq:LU}
  (1-2\epsilon)\ln\frac{1-2\epsilon}{x}+2\epsilon\ln\frac{2\epsilon}{1-x}
\end{equation}
(eq.~\texttt{LU}). It ends with ``$C_2$ is a numeric constant of order unity''
(line~500).

\section{The expectation $y$}

Line~446 states $y\ge1-2\epsilon$ ``as before'', but the earlier bound
$\langle\Psi_0,O_LO_R\Psi_0\rangle\ge1-2\epsilon$ is an expectation in the
ground state, while $y$ is an expectation in the product state $\sigma$. Since
$O_L$ and $O_R$ act on the two halves, $y=\langle O_L\rangle\langle O_R\rangle$
with expectations in $\Psi_0$. The bound that follows directly is weaker:
$\lVert O_BO_LO_R\Psi_0\rVert\ge1-\epsilon$ gives
$\lVert O_LO_R\Psi_0\rVert\ge1-\epsilon$, and for commuting positive
contractions $\langle O_L\rangle\ge\lVert O_L\Psi_0\rVert^2
\ge\lVert O_LO_R\Psi_0\rVert^2\ge(1-\epsilon)^2$, likewise for $O_R$, so
$y\ge(1-\epsilon)^4\ge1-4\epsilon$. We found no example in which the printed
bound fails. In a numerical search over two-qubit examples, with $O_B$ acting
on both qubits, the largest value of $(1-y)/\epsilon$ was about $0.98$. The
formalization uses the derived bound $1-4\epsilon$. It only moves the
threshold on $x$ below and the value of $C_2$, and leaves the coefficient
$l/\xi'$ of the strengthened doubling inequality unchanged.

With $\operatorname{tr}(O_BO_LO_R\sigma)\ge xy-\sqrt{x-x^2}\sqrt{y-y^2}$ and
the assumption of lines~484--490, which bounds the overlap term by $2\epsilon$,
this gives
\[
  x(1-4\epsilon)\le3\epsilon+2\sqrt{\epsilon x}.
\]
For $\epsilon\le1/68$, solving for $\sqrt x$ gives
$x\le(2+\sqrt{10})^2\epsilon\le36\,\epsilon$.

\section{The value of $C_2$}

The formalization proves~\eqref{eq:relent} with $C_2=7$, using only $2C_1>1$
and the nonnegativity of the mutual information. It splits on $\epsilon$.
\begin{itemize}
  \item For $\epsilon\le1/68$, $x\le36\epsilon$ and $1-2\epsilon\ge1/2$ give
    \[
      (1-2\epsilon)\ln\frac{1-2\epsilon}{x}
      \ge(1-2\epsilon)\ln\frac1\epsilon-\ln 72
      \ge(1-2\epsilon)\ln\frac1\epsilon-5 .
    \]
    The inequality $\ln u\ge1-1/u$ gives
    $2\epsilon\ln(2\epsilon/(1-x))\ge-1$. Since
    $\ln(1/\epsilon)=l/\xi'-\ln C_1$ and $\ln C_1>-\ln2>-1$,
    the bound~\eqref{eq:LU} is at least
    $(1-2\epsilon)l/\xi'-\ln C_1-7$.
  \item For $1/68<\epsilon<1/2$, the right side of~\eqref{eq:relent} equals
    $(1-2\epsilon)\ln(1/\epsilon)-2\epsilon\ln C_1-C_2$. This is at most
    $\ln68+\ln2-7<0$.
  \item For $\epsilon\ge1/2$, the right side is at most $-\ln C_1-7<0$.
\end{itemize}
In the last two cases the nonnegativity of the mutual information suffices.
The doubling iteration needs $\ln C_1+C_2\ge0$, which holds since
$\ln C_1>-1$.

\section{Formalization}

\leanid{TNLean.MPS.AreaLaw.Hastings.le_of_lindbladUhlmann} proves the case
analysis above for real numbers. Its hypotheses are $\Delta\ge0$ and, for
$\epsilon\le1/68$, a point $x\in(0,36\epsilon]$ at which~\eqref{eq:LU} is at
most $\Delta$. \leanid{TNLean.MPS.AreaLaw.Hastings.relentConst} is $C_2=7$.

\bibliographystyle{plain}
\bibliography{references}

\end{document}
